%=====================================================================
\chapter{The Entity-Relationship Model}\label{ch:ermodel}
%=====================================================================
\locator{2}{Part II --- Designing It}

\begin{outcomes}
By the end of this chapter you will be able to:
\begin{tight}
  \item \textbf{Distinguish} an entity, an entity type and an entity set, and
        \textbf{say} which of the three a table corresponds to.
  \item \textbf{Classify} an attribute as simple or composite, single- or
        multi-valued, stored or derived, and \textbf{explain} why each
        distinction changes what you build.
  \item \textbf{Identify} superkeys, candidate keys and the primary key of an
        entity type, and \textbf{justify} the choice between candidates.
  \item \textbf{Read and draw} a Chen ER diagram, and \textbf{state} the degree
        of a relationship.
  \item \textbf{Decide}, for a given noun in a requirements statement, whether
        it is an entity, an attribute or a relationship.
\end{tight}
\end{outcomes}

\begin{prereq}
\chref{whodoeswhat}: the designer's job --- deciding \emph{what is true}.
\chref{firstdatabase} for the \relation{}{} notation. No SQL is written in
this chapter; you will not touch \texttt{university.db} until
\chref{mapping}.
\end{prereq}

\begin{keyterms}
entity $\bullet$ entity type $\bullet$ entity set $\bullet$ attribute
$\bullet$ simple attribute $\bullet$ composite attribute $\bullet$
single-valued $\bullet$ multi-valued $\bullet$ stored attribute $\bullet$
derived attribute $\bullet$ null value $\bullet$ key attribute $\bullet$
superkey $\bullet$ candidate key $\bullet$ primary key $\bullet$ composite key
$\bullet$ surrogate key $\bullet$ relationship $\bullet$ relationship type
$\bullet$ degree $\bullet$ binary $\bullet$ ternary $\bullet$ recursive
relationship $\bullet$ role name $\bullet$ ER diagram $\bullet$ Chen notation
\end{keyterms}

\section{Design Happens on Paper}

You can create a table in eleven milliseconds (\chref{firstdatabase}). That is
exactly the problem: it is so cheap to make a table that people make the wrong
ones, and a wrong table is discovered in week fourteen when four programs
already read it.

The entity-relationship model is a way of arguing about a design before any of
it is expensive. It was proposed by Peter Chen in 1976 and it has outlived
every database product of that decade, for one reason: it is about the
\emph{university}, not about the database. Everything in this chapter could be
checked by the examinations officer, who has never heard of SQL.

\begin{alertbox}[HEC's outline puts this chapter after normalisation. This handout does not.]
The published course content lists \emph{normalization, functional
dependencies, normal forms} before \emph{entity relationship model, entity
sets, \dots, entity-relationship diagrams}. Taken literally, that normalises a
schema before designing one, which cannot be done --- normalisation
(\chref{normalforms}) takes an existing set of tables as its input.

\smallskip
All four of HEC's own recommended textbooks order it the other way. So does
this handout: design first (Part~II), map it to tables (Part~III), and only
then check the result for faults (Part~V). Every clause is still covered; the
Course Specification records the reordering and this is the one place it is
visible.
\end{alertbox}

\section{Entities: the Things}

\begin{definitionbox}[Entity, entity type, entity set]
An \term{entity} is one thing in the world that can be told apart from other
things: the student Ayesha Khan, the course CS-103, Tuesday's section of it.

\smallskip
An \term{entity type} is the \emph{description} of a kind of thing --- the
name \ent{student} together with the attributes every student has.

\smallskip
An \term{entity set} is all the entities of that type at one moment --- the
4{,}000 students enrolled today.

\smallskip
This is \chref{architecture}'s schema/instance distinction wearing different
words. The entity type is schema; the entity set is instance. A table
corresponds to the \textbf{entity type}; its rows are the entity set.
\end{definitionbox}

\begin{examplebox}[Three words people use for one thing, and when it matters]
\begin{tabular}{@{}L{3.6cm}L{5.0cm}L{5.4cm}@{}}
\textbf{Term} & \textbf{Example} & \textbf{Changes how often} \\
Entity        & Ayesha Khan, roll 2021-IT-001 & constantly \\
Entity set    & today's 4{,}000 students      & every admission and graduation \\
Entity type   & \ent{student}(roll\_no, name, programme, semester)
                                              & once a year, by argument \\
\end{tabular}

\smallskip
In conversation ``student'' means all three and no harm is done. In an
examination answer it must be the right one, and in a design meeting the
confusion is expensive: ``students change all the time'' is true of the set
and false of the type, and the two sentences lead to different databases.
\end{examplebox}

\section{Attributes, and the Four Distinctions That Matter}

An \term{attribute} is a property of an entity type. Students have a name and
a roll number. That much is obvious; the four ways attributes differ are not,
and each one changes what you have to build.

\begin{definitionbox}[Four distinctions]
\textbf{Simple or composite.} A \term{simple} (atomic) attribute cannot
usefully be divided: \texttt{semester}. A \term{composite} attribute has
meaningful parts: \texttt{address} = house, street, city, postcode.

\smallskip
\textbf{Single-valued or multi-valued.} \texttt{date\_of\_birth} has exactly
one value per student. \texttt{phone} may have several. A \term{multi-valued}
attribute is drawn with a double ellipse --- and it is the one kind that
\emph{cannot become a column}, which is \chref{mapping}'s problem.

\smallskip
\textbf{Stored or derived.} \texttt{date\_of\_birth} is \term{stored};
\texttt{age} is \term{derived} from it and today's date. A derived attribute
is drawn with a dashed ellipse and is normally \emph{not} stored at all.

\smallskip
\textbf{Key or non-key.} A \term{key attribute} has a different value for
every entity in the set. It is underlined.

\smallskip
A fifth property cuts across all four: an attribute may have no value for some
entity, which is the \term{null} of \chref{select} and a much deeper problem
than it looks.
\end{definitionbox}

\dsfig{%
\begin{tikzpicture}[
  ent/.style={draw=primary, line width=1.1pt, fill=primary!8,
              text=primary!60!black, font=\small\bfseries,
              minimum width=2.4cm, minimum height=0.85cm},
  att/.style={draw=secondary, line width=0.8pt, fill=white,
              text=secondary!60!black, font=\scriptsize, ellipse,
              inner sep=1.6pt, minimum width=1.5cm},
  keyatt/.style={att, text=primary, font=\scriptsize\bfseries},
  multi/.style={att, double, double distance=1.3pt, draw=purplehl,
                text=purplehl!70!black},
  deriv/.style={att, dashed, draw=neutral, text=neutral},
  lnk/.style={draw=secondary!50, line width=0.7pt},
  note/.style={font=\scriptsize\itshape, text=neutral}
]
  \node[ent] (s) at (0,0) {STUDENT};

  \node[keyatt] (k)  at (-3.9, 1.5) {\relkey{roll\_no}};
  \node[att]    (n)  at (-1.6, 2.1) {name};
  \node[att]    (pr) at ( 0.9, 2.1) {programme};
  \node[att]    (dob) at ( 3.3, 1.5) {date\_of\_birth};
  \node[deriv]  (ag) at ( 4.6, 0.1) {age};
  \node[multi]  (ph) at ( 3.1,-1.5) {phone};
  \node[att]    (ad) at (-1.4,-1.7) {address};
  \node[att]    (h)  at (-4.2,-0.8) {house};
  \node[att]    (st) at (-4.3,-1.8) {street};
  \node[att]    (ci) at (-3.4,-2.7) {city};

  \foreach \a in {k,n,pr,dob,ag,ph,ad}{\draw[lnk] (s) -- (\a);}
  \foreach \a in {h,st,ci}{\draw[lnk] (ad) -- (\a);}

  \node[note, anchor=west, text=purplehl] at (4.1,-2.2)
       {double = multi-valued};
  \node[note, anchor=west, text=neutral]  at (5.0,-0.5) {dashed = derived};
  \node[note, anchor=east, text=secondary!70!black] at (-4.6,-3.3)
       {composite: parts hang off the whole};
  \node[note, anchor=east, text=primary] at (-4.6,2.1) {underlined = key};
\end{tikzpicture}}%
{\ent{student} in Chen notation. Four of the five attribute kinds are
present. \texttt{phone} and \texttt{age} are the two that will not survive
contact with a table, and \chref{mapping} says what happens to each.}{erstudent}

\begin{conceptbox}[Why the distinctions are not decoration]
Each one has a consequence you will meet within three chapters.

\begin{tight}
  \item \textbf{Composite.} Store the parts, not the whole. A system that
        stores \texttt{address} as one text field can never answer ``how many
        students live in Bahawalpur'', and finding out costs a schema change
        after the data exists --- which is \chref{evolution}.
  \item \textbf{Multi-valued.} A column cannot hold three phone numbers. Not
        ``should not'': \emph{cannot}, if you want to search them. It becomes
        its own table, and that is a rule in \chref{mapping}.
  \item \textbf{Derived.} Storing \texttt{age} means storing a fact that
        becomes wrong overnight, silently, for about a quarter of the rows
        each year. Store the birth date and compute.
  \item \textbf{Key.} Everything else in the design hangs off this choice, and
        the next section is about how to make it.
\end{tight}
\end{conceptbox}

\section{Keys: Three Words, One Choice}

\begin{definitionbox}[Superkey, candidate key, primary key]
A \term{superkey} is any set of attributes whose values are different for
every entity. $\{$\texttt{roll\_no}$\}$ is one; so is
$\{$\texttt{roll\_no}, \texttt{name}$\}$, which is a superkey carrying a
passenger.

\smallskip
A \term{candidate key} is a superkey with no attribute to spare: remove any one
attribute and it stops being unique. $\{$\texttt{roll\_no}$\}$ is a candidate
key; $\{$\texttt{roll\_no}, \texttt{name}$\}$ is not.

\smallskip
The \term{primary key} is the one candidate key the designer picks to identify
entities throughout the database. The others are \term{alternate keys}.

\smallskip
A key made of more than one attribute is \term{composite}. A key invented by
the designer with no meaning in the world --- an auto-numbered id --- is a
\term{surrogate key}.
\end{definitionbox}

\begin{worked}[Finding every key of \ent{student}, and choosing one]
Attributes: \texttt{roll\_no}, \texttt{cnic}, \texttt{name},
\texttt{programme}, \texttt{semester}, \texttt{email}.

\medskip
\textbf{Step 1 --- what the university guarantees.} Not what today's data
happens to show. Ask, and write the answers down:
\begin{tight}
  \item Roll numbers are unique and never reissued. \emph{Guaranteed.}
  \item CNICs are unique. \emph{Guaranteed} --- but a first-semester student
        may not have one yet.
  \item Email addresses are issued one per student. \emph{Guaranteed} while
        enrolled; reused after five years.
  \item Names are not unique. Two students called Muhammad Ali is normal.
\end{tight}

\textbf{Step 2 --- candidate keys.} Three:
$\{$\texttt{roll\_no}$\}$, $\{$\texttt{cnic}$\}$, $\{$\texttt{email}$\}$.
$\{$\texttt{name}, \texttt{programme}, \texttt{semester}$\}$ is \emph{not} one
--- two BSIT second-semester students can share a name, and ``it has not
happened yet'' is not a guarantee.

\medskip
\textbf{Step 3 --- choose, on four grounds.}

\smallskip
\begin{tabular}{@{}L{2.2cm}L{2.4cm}L{2.4cm}L{2.2cm}L{3.0cm}@{}}
& \textbf{always present?} & \textbf{never changes?} & \textbf{short?}
& \textbf{verdict} \\
\texttt{roll\_no} & yes & yes & yes & \textbf{primary key} \\
\texttt{cnic}     & \textbf{no} & yes & yes & alternate, and nullable \\
\texttt{email}    & yes & \textbf{no} & no & alternate at best \\
\end{tabular}

\medskip
\textbf{The answer.} \relation{student}{\relkey{roll\_no}, cnic, name,
programme, semester, email}, with \texttt{cnic} and \texttt{email} declared
unique but not primary.

\smallskip
\textbf{What decided it.} Not uniqueness --- all three are unique. The primary
key must also be \emph{always present} and \emph{never changing}, because
every other table will refer to it (\chref{constraints}). \texttt{cnic} fails
the first test and \texttt{email} fails the second, and a primary key that can
change is a repair job spread across every table that referenced it.
\end{worked}

\begin{alertbox}[Three things students believe about keys, and what is wrong with each]
\begin{tight}
  \item \textbf{``The data is unique, so it is a key.''} A key is a
        \emph{constraint on all future data}, not an observation about today's.
        Check the current data by all means --- a duplicate proves it is not a
        key --- but no duplicate proves nothing.
  \item \textbf{``Add an auto-number id to every table and the problem goes
        away.''} The surrogate makes rows addressable; it does not make them
        \emph{distinguishable}. With an id and no other constraint you can
        insert the same student twice and the database is content. Surrogates
        are useful --- and whenever you use one, the real candidate key still
        has to be declared unique beside it.
  \item \textbf{``A composite key means bad design.''} It usually means the
        entity genuinely is identified by a combination. An enrolment is
        identified by \emph{which student} and \emph{which section}, and the
        pair is the honest key --- \chref{cardinality} and \chref{mapping}
        both rely on it.
\end{tight}
\end{alertbox}

\section{Relationships, and Degree}

\begin{definitionbox}[Relationship, relationship type, degree]
A \term{relationship} is an association among entities: Ayesha is enrolled in
Tuesday's section of CS-103.

\smallskip
A \term{relationship type} is the kind of association --- \ent{enrols},
between \ent{student} and \ent{section}. It is drawn as a diamond.

\smallskip
The \term{degree} of a relationship type is how many entity types take part:
\term{unary} (1), \term{binary} (2), \term{ternary} (3). Almost everything is
binary.

\smallskip
A \term{recursive} relationship is a unary one --- an entity type related to
itself, like a course being a prerequisite for another course. The two ends
need \term{role names} to tell them apart, because ``course'' appears twice.

\smallskip
A relationship type may have attributes of its own. \texttt{grade} belongs to
the \emph{enrolment}, not to the student and not to the section.
\end{definitionbox}

\dsfig{%
\begin{tikzpicture}[
  ent/.style={draw=primary, line width=1.1pt, fill=primary!8,
              text=primary!60!black, font=\scriptsize\bfseries,
              minimum width=2.0cm, minimum height=0.72cm},
  rel/.style={draw=greenhl, line width=1pt, fill=greenhl!8,
              text=greenhl!55!black, font=\scriptsize\bfseries,
              diamond, aspect=2.1, inner sep=1pt},
  att/.style={draw=secondary, line width=0.7pt, fill=white,
              text=secondary!60!black, font=\scriptsize, ellipse,
              inner sep=1.4pt},
  lnk/.style={draw=neutral!55, line width=0.7pt},
  rl/.style={font=\scriptsize\itshape, text=neutral},
  hdr/.style={font=\scriptsize\bfseries, text=primary!70!black}
]
  % --- binary
  \node[hdr] at (-4.6,1.9) {Binary};
  \node[ent] (st) at (-6.3,0.6) {STUDENT};
  \node[rel] (en) at (-4.0,0.6) {ENROLS};
  \node[ent] (se) at (-1.7,0.6) {SECTION};
  \node[att] (gr) at (-4.0,-0.9) {grade};
  \draw[lnk] (st) -- (en);  \draw[lnk] (en) -- (se);
  \draw[lnk] (en) -- (gr);
  \node[rl, text width=3.4cm, align=center] at (-4.0,-1.9)
       {the grade belongs to the \emph{relationship}, not to either end};

  % --- recursive
  \node[hdr] at (2.6,1.9) {Unary (recursive)};
  \node[ent] (co) at (2.6,0.6) {COURSE};
  \node[rel] (pr) at (2.6,-1.2) {REQUIRES};
  \draw[lnk] (co.west) .. controls (0.7,0.4) and (0.7,-1.1) .. (pr.west);
  \draw[lnk] (co.east) .. controls (4.5,0.4) and (4.5,-1.1) .. (pr.east);
  \node[rl, anchor=east] at (1.0,-0.3) {\emph{is prerequisite for}};
  \node[rl, anchor=west] at (4.2,-0.3) {\emph{depends on}};
  \node[rl, text width=3.8cm, align=center] at (2.6,-2.3)
       {one entity type, two ends: the role names are what tell them apart};
\end{tikzpicture}}%
{Two of the three degrees. \ent{enrols} carries an attribute of its own;
\ent{requires} joins \ent{course} to itself and is unreadable without
role names. The university database you are building contains both.}{erdegree}

\begin{examplebox}[Entity, attribute or relationship? The question that decides a design]
Every noun in a requirements document is one of the three, and getting it wrong
is the most common design error there is. Three tests, applied in order:

\begin{tight}
  \item \textbf{Does it have properties of its own?} If you can list two or
        three facts \emph{about} it, it is an entity.
  \item \textbf{Is there more than one of it per owner?} If yes, it is not an
        attribute.
  \item \textbf{Does it only exist because two other things met?} Then it is a
        relationship.
\end{tight}

\medskip
\begin{tabular}{@{}L{3.3cm}L{2.5cm}L{8.1cm}@{}}
\textbf{Noun} & \textbf{Verdict} & \textbf{Why} \\
roll number   & attribute    & no properties of its own \\
department    & \textbf{entity} & it has a name, a head, a building --- and
                                 fifty courses point at it \\
grade         & \textbf{relationship attribute} & exists only because a
                                 student met a section \\
prerequisite  & \textbf{relationship} & joins course to course \\
semester      & it depends  & an attribute of \ent{student} if it means
                              ``which year they are in''; an \emph{entity} if
                              the university runs Fall 2026 and Spring 2027 as
                              distinct things with dates \\
\end{tabular}

\smallskip
The last row is the honest one. The same word is an attribute in one
university and an entity in another, and no rule decides it --- the
\emph{requirement} does. That is why design is a conversation with the
examinations officer and not a procedure.
\end{examplebox}

\begin{pitfall}
\begin{tight}
  \item \textbf{Making an entity out of every noun.} ``Name'' is not an
        entity. Apply the three tests.
  \item \textbf{Hanging a relationship's attribute on an entity.} Putting
        \texttt{grade} on \ent{student} gives each student one grade for
        the whole degree. It belongs on \ent{enrols}.
  \item \textbf{Choosing a primary key that can change.} Email, phone, name.
        Every table that referenced it has to be repaired with it.
  \item \textbf{Treating today's data as a guarantee.} No duplicate names today
        does not make \texttt{name} a key.
  \item \textbf{Drawing a recursive relationship without role names.} The
        diagram is then genuinely ambiguous --- nobody can tell which end is
        the prerequisite.
  \item \textbf{Storing a derived attribute because it is convenient.}
        \texttt{age} is wrong for a quarter of your rows within a year, and
        nothing reports it.
\end{tight}
\end{pitfall}

\begin{lab}[Turn a page of English into an ER diagram]
No database this week. Paper, and one argument.

\begin{lstlisting}
# --- 1. The requirement, as the examinations officer states it -------

"Every student has a roll number, a name, a CNIC, one or more phone
 numbers, a postal address, and a date of birth. Students belong to
 one programme. A course has a code, a title and a number of credit
 hours, and belongs to a department. A course may require other
 courses as prerequisites. Each semester a course runs as one or
 more sections; a section has a number, a room and a timetable slot,
 and is taught by one instructor. Students enrol in sections and
 receive a grade. We need to know each student's age, and how many
 students are in each section."
\end{lstlisting}

\begin{lstlisting}
# --- 2. List every noun --------------------------------------------
    Write them in one column. There are about twenty.

# --- 3. Classify each one ------------------------------------------
    entity / attribute / relationship, using the three tests.
    Write ONE SENTENCE of justification for each. The sentence is
    the exercise; the label is not.

# --- 4. Mark the attribute kinds -----------------------------------
    composite, multi-valued, derived, key. There is at least one
    of each in the paragraph above.

# --- 5. Find the candidate keys ------------------------------------
    For STUDENT, COURSE and SECTION. For each, say what the
    university guarantees -- not what sample data shows.

# --- 6. Draw it, Chen notation -------------------------------------
    Rectangles, ellipses, diamonds. Role names on the recursive
    relationship. Do NOT draw cardinalities yet; that is Chapter 6.

# --- 7. Find the two traps -----------------------------------------
    Two things in the paragraph must NOT become columns, and one
    must not be stored at all. Name all three and say why.

# --- 8. Hand it to somebody else -----------------------------------
    Give your diagram, without the paragraph, to another student.
    Ask them to read it back as English. Every sentence they get
    wrong is a defect in your diagram, not in their reading.
\end{lstlisting}

\smallskip
\textbf{What to notice.} Step 8 is the one that matters and the one people
skip. An ER diagram's whole purpose is to be checked by somebody who cannot
read SQL. If it cannot be read back, it has failed at the only job it has ---
and it is still cheap to fix, which it will not be in \chref{mapping}.

\smallskip
\textbf{Keep this diagram.} Chapter~6 adds cardinalities to it, \chref{mapping}
turns it into tables, and it is part of your project submission.
\end{lab}

\begin{checkpoint}
\begin{tightnum}
  \item A colleague proposes \texttt{phone} as a column of \ent{student},
        because ``almost every student has one number''. Give the argument
        against in three sentences, and say what you would do instead and what
        it costs.
  \item \ent{instructor} has \texttt{emp\_id}, \texttt{cnic}, \texttt{name},
        \texttt{email}, \texttt{department}, \texttt{office\_no}. List every
        candidate key you can justify, state the guarantee each one rests on,
        choose a primary key and defend the choice on all four grounds in the
        worked example.
  \item Explain, to somebody who has written programs but never designed a
        database, why \texttt{grade} belongs to the \emph{relationship} between
        a student and a section rather than to either of them. Then say what
        goes wrong if it is put on \ent{student}, and what goes wrong if it
        is put on \ent{section}.
\end{tightnum}
\end{checkpoint}

\begin{chaptersummary}
\begin{tight}
  \item An \textbf{entity} is one thing; an \textbf{entity type} describes a
        kind of thing; an \textbf{entity set} is all of them today. Type is
        schema, set is instance, and a table corresponds to the type.
  \item Attributes differ in four ways: simple/\textbf{composite},
        single-/\textbf{multi-valued}, stored/\textbf{derived}, key/non-key ---
        and each distinction has a consequence in Part~III.
  \item A \textbf{multi-valued} attribute cannot become a column; a
        \textbf{derived} one should not be stored; a \textbf{composite} one is
        stored as its parts.
  \item A \textbf{superkey} is unique; a \textbf{candidate key} is unique with
        nothing to spare; the \textbf{primary key} is the candidate the
        designer picks, and the others are alternate keys.
  \item Choose a primary key on four grounds: unique, \textbf{always present},
        \textbf{never changing}, short. Uniqueness alone decides nothing when
        there are three candidates.
  \item A key is a constraint on \emph{all future data}. A duplicate disproves
        a key; the absence of one proves nothing.
  \item A \textbf{relationship} associates entities; its \textbf{degree} is how
        many entity types take part. A \textbf{recursive} relationship joins a
        type to itself and needs \textbf{role names} to be readable.
  \item A relationship can carry attributes of its own. \texttt{grade} belongs
        to the enrolment, not to the student and not to the section.
  \item Every noun is an entity, an attribute or a relationship. The same word
        can be different things in different universities; the requirement
        decides, not a rule.
  \item HEC's clause list puts normalisation before the ER model. This handout
        reverses the two, because normalisation takes tables as its input and
        there are none until Part~III.
\end{tight}
\end{chaptersummary}

\begin{reviewq}
  \item \textbf{(MCQ)} A double ellipse in Chen notation marks an attribute
        that is: (a)~derived; (b)~composite; (c)~multi-valued; (d)~a key.
  \item \textbf{(MCQ)} $\{$\texttt{roll\_no}, \texttt{name}$\}$, where
        \texttt{roll\_no} is unique, is: (a)~a candidate key; (b)~a superkey
        but not a candidate key; (c)~neither; (d)~the primary key.
  \item \textbf{(Short)} Give an example of a derived attribute other than age,
        and say what goes wrong if it is stored.
  \item \textbf{(Short)} What is the degree of the \ent{requires}
        relationship between courses, and why does it need role names when a
        binary relationship usually does not?
  \item \textbf{(Short)} A surrogate key is added to a table that already has a
        natural candidate key. What must still be declared, and what breaks if
        it is not?
  \item \textbf{(Applied)} A library wants to record books, copies of books,
        members and loans. Decide which of the four are entity types, identify
        the relationships and their degrees, choose a primary key for each
        entity type with its justification, and say which attributes of a loan
        belong to the relationship rather than to either end.
  \item \textbf{(Applied)} Take the \ent{student} diagram in
        Figure~\ref{fig:erstudent} and the requirement ``we must be able to
        list every student in a given city, and ring any of a student's
        numbers''. State precisely which parts of the diagram now cannot be
        implemented as columns, what you would build instead, and what the
        design would have cost if this requirement had arrived after the
        database was full.
\end{reviewq}
